\subsection{Discrete Fourier Transforms and Circulant Matrices}

\begin{definition} \label{def:discrete-fourier-matrix} \hlt{Discrete Fourier Matrix}\\
For $n\geq1$, set $\omega_n=e^{2\pi i/n}$ and $F_n\defeq[\omega_n^{jk}]_{j,k=0}^{n-1}$.\\
Then $F_n\mathbf{c}$ is the unnormalised DFT of $\mathbf{c}\in\C^n$; $Q_n\defeq n^{-1/2}F_n$ is the normalized Fourier matrix.
\end{definition}

\begin{definition} \label{def:circulant-matrix-cyclic-shift} \hlt{Circulant Matrix}\\
On the $0$-indexed basis of $\C^n$, define the \hlt{cyclic shift}
$P_n\mathbf{e}_j=\mathbf{e}_{(j-1)\operatorname{mod} n}$.\\
A matrix $C\in\operatorname{Mat}(n,n,\C)$ is \hlt{circulant} if, for some $a_0,\ldots,a_{n-1}\in\C$,
\begin{align}
C&=a_0I_n+a_1P_n+\cdots+a_{n-1}P_n^{n-1}=p(P_n), \nonumber\\
p(z)&=a_0+a_1z+\cdots+a_{n-1}z^{n-1} \nonumber
\end{align}
\end{definition}

\begin{theorem} \label{thm:fourier-basis-diagonalizes-cyclic-circulants} \hlt{Fourier Basis for Cyclic Shifts and Circulants}\\
Let $n\geq1$, let $\omega_n\defeq e^{2\pi i/n}$, and let $P_n$ be the cyclic shift.\\
For $0\leq k<n$, set $\mathbf{f}_k\defeq(\omega_n^{jk})_{j=0}^{n-1}$; $F_n\defeq[\mathbf{f}_0\ \cdots\ \mathbf{f}_{n-1}]$; and $\Omega_n\defeq\operatorname{diag}(1,\omega_n,\ldots,\omega_n^{n-1})$.\\
Then $F_n^*F_n=nI_n$, so $Q_n\defeq n^{-1/2}F_n$ is unitary, and $P_nF_n=F_n\Omega_n$.\\
If $C=p(P_n)$ is circulant, then $Q_n^*CQ_n=p(\Omega_n) =\operatorname{diag}\bigl(p(1),p(\omega_n),\ldots,p(\omega_n^{n-1})\bigr)$.
\begin{proof}
For $0\leq k,\ell\leq n-1$, $\innerprod{\mathbf{f}_k,\mathbf{f}_\ell} =\sum\limits_{r=0}^{n-1}(\omega_n^{k-\ell})^r =n\delta_{k\ell}$.
For $k=\ell$ the sum is $n$; for $k\neq\ell$, $W=\omega_n^{k-\ell}$ has $W\neq1$ and $W^n=1$, so $(1-W)\sum\limits_{r=0}^{n-1}W^r=0$. Hence $F_n^*F_n=nI_n$, and Theorem \ref{thm:unitary-matrix-equivalent-conditions} makes $Q_n$ unitary.\\
Also $P_n\mathbf{f}_k=\omega_n^k\mathbf{f}_k$, so $Q_n^*P_nQ_n=\Omega_n$; applying $p$ gives the formula.\\
\hlt{Structural Commentary}: Fourier coordinates diagonalize every circulant.\\
\hlt{Key Insight}: Orthogonality of roots of unity turns a cyclic matrix problem into scalar evaluation.
\end{proof}
\end{theorem}

\begin{figure}[H]
  \centering
  \begin{tikzpicture}[
  >=Stealth,
  x=0.92cm,
  y=0.9cm,
  every node/.style={font=\small},
  box/.style={draw=gray!55, rounded corners=3pt, fill=gray!6, align=center, minimum width=2.15cm, minimum height=0.8cm},
  result/.style={box, fill=green!7},
  arrow/.style={->, line width=0.85pt, blue!70!black},
  lbl/.style={fill=white, text=gray!70!black, inner xsep=2pt, inner ysep=1pt}
]
  \node[box] (shift) at (0,1.1) {$P_n$};
  \node[box] (diagshift) at (4.15,1.1) {$\Omega_n$};
  \node[box] (circ) at (0,-1.1) {$C=p(P_n)$};
  \node[result] (diagcirc) at (4.15,-1.1) {$p(\Omega_n)$};

  \draw[arrow] (shift) -- node[lbl, above] {$Q_n^*(\,\cdot\,)Q_n$} (diagshift);
  \draw[arrow] (circ) -- node[lbl, below] {$Q_n^*(\,\cdot\,)Q_n$} (diagcirc);
  \draw[arrow] (shift) -- node[lbl, left] {$p(\cdot)$} (circ);
  \draw[arrow] (diagshift) -- node[lbl, right] {$p(\cdot)$} (diagcirc);
\end{tikzpicture}

  \caption{Fourier diagonalizes $P_n$; polynomial evaluation diagonalizes $C=p(P_n)$.}
  \label{fig:circulant-fourier-diagonalization}
\end{figure}

\begin{proposition} \label{prop:fourier-matrix-inverse} \hlt{Inverse of the Fourier Matrix}\\
The Fourier matrix satisfies $F_n^*F_n=nI_n$, and therefore $F_n^{-1}=\frac1nF_n^*$.
\begin{proof}
Theorem \ref{thm:fourier-basis-diagonalizes-cyclic-circulants} gives $Q_n^*Q_n=I_n$ for $Q_n=n^{-1/2}F_n$. Multiplying by $n$ gives $F_n^*F_n=nI_n$, so $n^{-1}F_n^*$ is a left inverse of $F_n$. Since $F_n$ is square, Theorem \ref{thm:square-one-sided-inverse-is-two-sided} makes it the inverse.\\
\hlt{Structural Commentary}: Fourier inversion is conjugate transpose plus one scalar normalization.\\
\hlt{Key Insight}: The inverse is simple because the Fourier columns are orthogonal.
\end{proof}
\end{proposition}
